\documentclass[12pt,letterpaper]{article}
\usepackage[utf8]{inputenc}

% ── Packages ──────────────────────────────────────────────────────────────────
\usepackage[margin=1in]{geometry}
\usepackage{booktabs}
\usepackage{longtable}
\usepackage{array}
\usepackage{xcolor}
\usepackage{colortbl}
\usepackage{siunitx}
\usepackage{amsmath}
\usepackage{hyperref}
\usepackage{graphicx}
\usepackage{enumitem}
\usepackage{fancyhdr}
\usepackage{titlesec}
\usepackage{caption}
\usepackage{tikz}
\usepackage[american, siunitx]{circuitikz}
\usepackage{float}

\usepackage{parskip}

% ── Colors ────────────────────────────────────────────────────────────────────
\definecolor{violet}{RGB}{109,40,217}
\definecolor{ledblue}{RGB}{29,111,181}
\definecolor{ledgreen}{RGB}{21,128,61}
\definecolor{ledred}{RGB}{185,28,28}
\definecolor{rowgray}{RGB}{245,245,245}
\definecolor{headblue}{RGB}{30,64,175}
\definecolor{warnred}{RGB}{185,28,28}
\definecolor{notegreen}{RGB}{21,128,61}
\definecolor{lightviolet}{RGB}{237,233,254}
\definecolor{lightblue}{RGB}{219,234,254}
\definecolor{lightgreen}{RGB}{220,252,231}
\definecolor{lightred}{RGB}{254,226,226}

% ── Header / Footer ───────────────────────────────────────────────────────────
\pagestyle{fancy}
\fancyhf{}
\rhead{\small Color Analyzer LED Driver Design}
\lhead{\small \leftmark}
\rfoot{\small Page \thepage}
%\cfoot{\small main\_board\_design.tex}
\cfoot{\small \jobname}
\lfoot{\small \today}

% ── Section formatting ────────────────────────────────────────────────────────
\titleformat{\section}{\large\bfseries\color{headblue}}{{\thesection}}{1em}{}[\titlerule]
\titleformat{\subsection}{\normalsize\bfseries}{{\thesubsection}}{1em}{}

% ── Hyperref setup ────────────────────────────────────────────────────────────
\hypersetup{colorlinks=true, linkcolor=headblue, urlcolor=headblue}

% ── siunitx setup ─────────────────────────────────────────────────────────────
\sisetup{per-mode=symbol}

% ── Column types ─────────────────────────────────────────────────────────────
\newcolumntype{C}[1]{>{\centering\arraybackslash}p{#1}}
\newcolumntype{L}[1]{>{\raggedright\arraybackslash}p{#1}}
\newcolumntype{R}[1]{>{\raggedleft\arraybackslash}p{#1}}

% ─────────────────────────────────────────────────────────────────────────────
\begin{document}
	
	\begin{titlepage}
		\centering
		\vspace*{2cm}
		{\Huge\bfseries\color{headblue} Main Board Design\par}
		\vspace{0.5cm}
		{\Large Circuit Analyses and Component Selection\par}
		\vspace{2cm}
		{\large Raspberry Pi Zero W \& Zero 2 W\par}
		\vspace{0.5cm}
		{\large Adafruit TFT 2.8in Touchscreen Display\\TSC2007 -- SPI and I2C\par}
		\vspace{0.5cm}
		{\large LED boards -- Testing -- Connectors -- Power\par}
		\vspace{3cm}
		{\normalsize\today\par}
		\vfill
		{\small This document covers design ideas, component selection, and sub-circuit diagrams.\\
			It is written for a technically-minded reader who may not have prior electronics design experience.}
	\end{titlepage}
	
	\tableofcontents
	
	
	% ── Section files ─────────────────────────────────────────────────────────
	% All .tex files must be in the same folder as this parent document.
	% Use \include{filename} if you want each section to start on a fresh page.
	% Use \input{filename} so that no \clearpage is forced between sections.
	
	%	\documentclass[border=10pt]{standalone}
% 'standalone' produces a tight bounding box around the drawing.
% 'border=10pt' adds a small margin around the content.

\usepackage{circuitikz}
% The main package for drawing electronic circuits in LaTeX.

\usepackage{xcolor}
% Lets us name colors for annotation text.

% ── Convenience color for annotation comments ──────────────
\definecolor{annotate}{RGB}{0,100,180}

\begin{document}


\section{Magnetic Buzzer}

\subsection{Overview}

The alert sounder is a passive magnetic buzzer driven by a pulse-width modulated (PWM)
signal from the Raspberry Pi Zero. A passive magnetic buzzer is preferred over an active
buzzer or piezoelectric transducer because its tone is warm and rounded rather than harsh,
and because both pitch and volume are controllable in software via the PWM frequency and
duty cycle respectively.

The buzzer operates from the \SI{5}{\volt} supply rail and draws approximately
\SI{30}{\milli\ampere} to \SI{40}{\milli\ampere} at full drive. Since this exceeds the
safe output current of a GPIO pin, and because the supply voltage exceeds the
\SI{3.3}{\volt} GPIO logic level, a small-signal NPN bipolar junction transistor is
used as a low-side switch between the buzzer and ground.

\subsection{Circuit Description}

The GPIO pin drives the base of the NPN transistor through a current-limiting resistor
R1. When the GPIO is high, the transistor saturates and connects the buzzer's negative
terminal to ground, completing the circuit through the \SI{5}{\volt} supply. When the
GPIO is low, the transistor is cut off and the buzzer is de-energised.

Because the buzzer is an inductive load, a flyback diode is placed across its terminals
to suppress the reverse voltage spike that occurs when the transistor switches off.
Without this diode the spike could damage the transistor.

Volume is controlled by adjusting the PWM duty cycle. A duty cycle of approximately
50\% produces full volume; reducing the duty cycle lowers the average power delivered
to the buzzer and therefore reduces loudness. Pitch is controlled by adjusting the PWM
frequency. The range of \SI{2000}{\hertz} to \SI{2500}{\hertz} produces a clear,
pleasant tone that is audible in both laboratory and outdoor environments. The buzzer
is rated at approximately \SI{85}{\decibel} to \SI{90}{\decibel} at \SI{10}{\centi\meter}.

The complete circuit is shown below.

\begin{verbatim}
    RPi GPIO (PWM) ---[ R1 1k ]---+--- MMBT3904 Base
                                  |
                             MMBT3904 Emitter --- GND

    +5V ---[ Buzzer CEM-1203(42) ]---+--- MMBT3904 Collector
                                     |
                               1N4148W (flyback)
                               Anode  -> Collector side
                               Cathode -> +5V side
\end{verbatim}

\subsection{Parts List}

\begin{table}[h]
\centering
\begin{tabular}{llll}
\hline
\textbf{Reference} & \textbf{Part Number} & \textbf{Description} & \textbf{Package} \\
\hline
BZ1 & CEM-1203(42)  & Passive magnetic buzzer, \SI{5}{\volt}, \SI{2048}{\hertz} rated & 12\,mm through-hole \\
D1  & 1N4148W       & Fast switching diode, flyback protection                        & SOD-123 \\
Q1  & MMBT3904      & NPN small-signal BJT, \SI{40}{\volt}, \SI{200}{\milli\ampere}  & SOT-23  \\
R1  & R1K-0805      & Resistor \SI{1}{\kilo\ohm}, base current limiter               & 0805    \\
\hline
\end{tabular}
\caption{Magnetic buzzer circuit components}
\label{tab:buzzer_parts}
\end{table}

\subsection{Programming Notes}

The buzzer is controlled using the RPi.GPIO software PWM facility. Any GPIO pin may
be used, but a hardware PWM pin (GPIO 18 or GPIO 12 on the Raspberry Pi Zero) is
preferred because software PWM can exhibit jitter under system load, which may cause
audible roughness in the tone.

Key relationships are:

\begin{itemize}
    \item \textbf{Frequency} sets pitch. The range \SI{1500}{\hertz} to
          \SI{3000}{\hertz} is recommended. Values below \SI{1000}{\hertz} may be too
          quiet for some magnetic buzzers, and values above \SI{4000}{\hertz} become
          increasingly harsh.
    \item \textbf{Duty cycle} sets volume. A value of 50 (50\%) gives maximum loudness.
          Values below approximately 10 may be too quiet to be useful. Values above 50
          do not increase loudness and may introduce distortion.
    \item \textbf{Duration} is implemented as a blocking \texttt{time.sleep()} call.
          Interruption logic is not implemented in this version.
\end{itemize}

The GPIO pin must be initialised as an output before the PWM object is created. The
PWM object should be started with a duty cycle of zero and frequency set before each
call to avoid a brief click at startup.

\subsection{Example Code}

\begin{verbatim}
import RPi.GPIO as GPIO
import time

BUZZER_PIN = 18  # Hardware PWM pin recommended

GPIO.setmode(GPIO.BCM)
GPIO.setup(BUZZER_PIN, GPIO.OUT)

_pwm = GPIO.PWM(BUZZER_PIN, 2000)
_pwm.start(0)


def sound_alert(frequency: int = 2000,
                volume: int = 50,
                duration: float = 1.0) -> None:
    """
    Sound the magnetic buzzer for a fixed duration.

    Parameters
    ----------
    frequency : int
        Tone frequency in Hz. Recommended range 1500-3000.
        Default 2000 Hz.
    volume : int
        Duty cycle as a percentage, 0-50. Values above 50
        do not increase loudness. Default 50 (full volume).
    duration : float
        Duration in seconds to sound the alert. Blocking.
        Default 1.0 second.
    """
    volume = max(0, min(50, volume))       # clamp to 0-50
    frequency = max(100, min(10000, frequency))  # clamp to safe range

    _pwm.ChangeFrequency(frequency)
    _pwm.ChangeDutyCycle(volume)
    time.sleep(duration)
    _pwm.ChangeDutyCycle(0)                # silence buzzer, leave PWM running
\end{verbatim}

\noindent To sound a two-second alert at \SI{2500}{\hertz} at three-quarter volume:

\begin{verbatim}
    sound_alert(frequency=2500, volume=37, duration=2.0)
\end{verbatim}

\noindent To clean up GPIO resources on program exit:

\begin{verbatim}
    _pwm.stop()
    GPIO.cleanup()
\end{verbatim}

	%	\section{Extending MOSFETs with Shift Register}

\subsection{Overview}

The Raspberry Pi Zero provides a limited number of GPIO pins. Where multiple independent
loads must be switched, a serial-to-parallel shift register allows eight MOSFET gates to
be controlled from three GPIO pins. The device chosen is the \textbf{74HC595}, an
8-bit serial-in, parallel-out shift register that operates from \SI{3.3}{\volt} and is
widely available in an SMD package.

Data is clocked into the shift register serially and latched to the eight parallel outputs
simultaneously, ensuring that all outputs update at the same instant. Although the
74HC595 can be driven by bit-banging three GPIO pins, the implementation here uses the
hardware \textbf{SPI1} peripheral on the Raspberry Pi Zero for reliable timing and
minimal CPU overhead. Only one MOSFET is energised at a time; the controlling byte
therefore always has exactly one bit set.

\subsection{Circuit Description}

The Raspberry Pi Zero SPI1 peripheral provides three signals: MOSI (data), SCLK (clock),
and CE (chip enable, used here as the latch signal). Each signal passes through a
\SI{470}{\ohm} series resistor before reaching the 74HC595 to protect the GPIO pins
from transient currents.

The 74HC595 output enable pin (\textoverline{OE}) is tied permanently to ground so
that the parallel outputs are always active. The master reset pin (\textoverline{MR})
is tied permanently to \SI{3.3}{\volt} to prevent accidental clearing of the register.
A \SI{100}{\nano\farad} bypass capacitor is placed directly across the Vcc and GND
pins of the IC.

Each of the eight parallel outputs drives a MOSFET gate through a \SI{500}{\ohm}
series resistor. A \SI{10}{\kilo\ohm} resistor from gate to ground ensures the gate
is pulled firmly low when the corresponding output is de-asserted, preventing the
MOSFET from floating into an undefined state.

\subsection{Circuit Diagram}

\begin{verbatim}
  Raspberry Pi Zero                  74HC595 (SOIC-16)
  (SPI1)                             +----------------+
                          3.3V -----| Vcc(16) (8)GND |------ GND
                                    |                |
   GPIO 20 (MOSI) --[ 470R ]--------| DS   (pin 14)  |   [SER]
   GPIO 21 (SCLK) --[ 470R ]--------| SHCP (pin 11)  |   [SRCLK]
   GPIO 18 (CE0)  --[ 470R ]--------| STCP (pin 12)  |   [RCLK]
                                    |                |
                            GND ----| /OE  (pin 13)  |  (always enabled)
                           3.3V ----| /MR  (pin 10)  |  (master reset off)[SRCLR]
                                    |                |
                                    |   Q0  (pin 15) |--+
                                    |   Q1  (pin 1)  |--+  (Q1..Q7 identical)
                                    |   Q2  (pin 2)  |--+
                                    |   Q3  (pin 3)  |--+
                                    |   Q4  (pin 4)  |--+
                                    |   Q5  (pin 5)  |--+
                                    |   Q6  (pin 6)  |--+
                                    |   Q7  (pin 7)  |--+
                                    |       (pin 9)  |--x  (Q7S serial data out)
                          100nF     |                |
                                    +----------------+
  --||-- Vcc to GND, close to IC
 
  One MOSFET channel (repeated x8):

  74HC595 Qn ---[ 500R ]---+--- MOSFET Gate (AO3400A or similar)
                            |
                          10k
                            |
                           GND                MOSFET Source --- GND
                                              MOSFET Drain  --- Load (-)
\end{verbatim}

\subsection{Parts List}

\begin{table}[h]
\centering
\begin{tabular}{llll}
\hline
\textbf{Reference} & \textbf{Part Number} & \textbf{Description} & \textbf{Package} \\
\hline
U1      & 74HC595           & 8-bit shift register, serial-in parallel-out  & SOIC-16 \\
Q1--Q8  & AO3400A           & N-channel MOSFET, \SI{30}{\volt}, \SI{5.7}{\ampere} & SOT-23 \\
R1--R3  & R470-0603         & \SI{470}{\ohm} GPIO series resistors          & 0603 \\
R4--R11 & R500-0603         & \SI{500}{\ohm} MOSFET gate resistors          & 0603 \\
R12--R19 & R10K-0603        & \SI{10}{\kilo\ohm} gate pull-down resistors   & 0603 \\
C1      & 100nF ceramic     & Bypass capacitor, Vcc to GND                  & 0603 \\
\hline
\end{tabular}
\caption{Shift register MOSFET controller components}
\label{tab:shiftreg_parts}
\end{table}

\subsection{Programming Notes}

The SPI1 peripheral on the Raspberry Pi Zero is accessed via the \texttt{spidev} library.
SPI mode 0 (CPOL=0, CPHA=0) is compatible with the 74HC595. A clock speed of
\SI{500}{\kilo\hertz} is more than adequate; the 74HC595 is rated to several megahertz
but there is no benefit to running faster for low-speed MOSFET switching.

A single byte is transferred per operation. The byte has exactly one bit set to select
one MOSFET channel, or is zero to de-energise all channels. The \texttt{xfer2()} call
clocks out the data and automatically toggles the chip enable line, which is connected
to the latch (STCP) pin, causing the outputs to update at the end of the transfer.

Note that SPI1 on the Raspberry Pi Zero requires the overlay \texttt{spi1-1cs} to be
enabled in \texttt{/boot/config.txt}:

\begin{verbatim}
    dtoverlay=spi1-1cs
\end{verbatim}

\subsection{Example Code}

\begin{verbatim}
import spidev

SPI_BUS = 1       # SPI1
SPI_DEVICE = 0    # CE0

spi = spidev.SpiDev()
spi.open(SPI_BUS, SPI_DEVICE)
spi.max_speed_hz = 500000
spi.mode = 0


def set_mosfet(channel: int) -> None:
    """
    Energise a single MOSFET by channel number.

    Parameters
    ----------
    channel : int
        Channel to activate, 0 through 7.
        Pass -1 or any out-of-range value to de-energise all channels.
    """
    if 0 <= channel <= 7:
        byte = 1 << channel
    else:
        byte = 0
    spi.xfer2([byte])


def all_off() -> None:
    """De-energise all MOSFET channels."""
    spi.xfer2([0x00])
\end{verbatim}

\noindent To activate channel 4 then turn everything off:

\begin{verbatim}
    set_mosfet(4)
    time.sleep(1.0)
    all_off()
\end{verbatim}


% ---------------------------------------------------------------------------
\section{High-Side Control of 5\,V Load with MOSFET}

\subsection{Overview}

In circuits requiring both high-side and low-side switching, the load is interrupted at
both the positive supply terminal and the ground return. The low-side switch is an
N-channel MOSFET (AO3400A), which is straightforward to drive from a \SI{3.3}{\volt}
GPIO. The high-side switch requires a P-channel MOSFET because its source terminal
connects to the positive supply rail.

The P-channel MOSFET chosen is the \textbf{AO3401A}, the complementary P-channel
counterpart to the AO3400A. Both devices share the SOT-23 package and are produced
by the same manufacturer, making them a well-matched pair. The AO3401A is rated at
\SI{-30}{\volt} Vds and \SI{-4}{\ampere} continuous drain current, well in excess of
the anticipated load current.

\subsection{The Level-Shifting Problem}

A P-channel MOSFET turns on when its gate is pulled sufficiently below its source
voltage. With the source at \SI{+5}{\volt}, the gate must be pulled close to ground
to turn the device on, and must be pulled to \SI{+5}{\volt} to turn it off. A
\SI{3.3}{\volt} GPIO cannot swing to \SI{+5}{\volt} and therefore cannot turn the
device fully off. A small-signal NPN transistor is used to solve this: it pulls the
gate to ground when the GPIO is high (turning the P-channel on), and allows the gate
to be pulled to \SI{+5}{\volt} through a resistor when the GPIO is low (turning the
P-channel off).

\subsection{Circuit Description}

The GPIO drives the base of an MMBT3904 NPN transistor through a \SI{10}{\kilo\ohm}
base resistor. The collector of the NPN connects to the gate of the AO3401A. A
\SI{10}{\kilo\ohm} resistor from the gate to \SI{+5}{\volt} pulls the gate high
when the NPN is off, ensuring the P-channel is firmly off. When the NPN saturates,
it pulls the gate to ground, turning the P-channel on.

The logic is inverted by the NPN stage: GPIO high turns the high-side MOSFET on,
and GPIO low turns it off. This matches the behaviour of the low-side N-channel
switch, so both MOSFETs are enabled by a high logic level on their respective GPIO
pins. The load is energised only when both GPIO pins are high simultaneously.

When the supply voltage is \SI{5}{\volt} and the gate is pulled to ground by the NPN,
Vgs of the AO3401A is \SI{0}{\volt} $-$ \SI{5}{\volt} $=$ \SI{-5}{\volt}. The
AO3401A gate threshold is \SI{-1}{\volt}, so the device is driven well into
conduction.

\subsection{Circuit Diagram}

\begin{verbatim}
  High-side control:

  RPi GPIO A --[ 10k ]--+-- MMBT3904 Base
                         |
                    MMBT3904 Emitter ----------- GND
                    MMBT3904 Collector --+------ AO3401A Gate
                                        |
                                       10k
                                        |
                                      +5V

  +5V ----------- AO3401A Source
                  AO3401A Drain ----+---- Load (+)
                                    |
                                  Load (-)
                                    |
                               AO3400A Drain
                               AO3400A Source --- GND
                               AO3400A Gate --[ 500R ]--+-- RPi GPIO B
                                                        |
                                                       10k
                                                        |
                                                       GND

  Load energised when: GPIO A = HIGH and GPIO B = HIGH
\end{verbatim}

\subsection{Parts List}

\begin{table}[h]
\centering
\begin{tabular}{llll}
\hline
\textbf{Reference} & \textbf{Part Number} & \textbf{Description} & \textbf{Package} \\
\hline
Q1 & AO3401A   & P-channel MOSFET, \SI{-30}{\volt}, \SI{-4}{\ampere}, high-side switch & SOT-23 \\
Q2 & AO3400A   & N-channel MOSFET, \SI{30}{\volt}, \SI{5.7}{\ampere}, low-side switch  & SOT-23 \\
Q3 & MMBT3904  & NPN BJT, \SI{40}{\volt}, \SI{200}{\milli\ampere}, gate driver         & SOT-23 \\
R1 & R10K-0603 & \SI{10}{\kilo\ohm} NPN base resistor                                  & 0603   \\
R2 & R10K-0603 & \SI{10}{\kilo\ohm} P-channel gate pull-up to \SI{5}{\volt}            & 0603   \\
R3 & R500-0603 & \SI{500}{\ohm} N-channel gate series resistor                         & 0603   \\
R4 & R10K-0603 & \SI{10}{\kilo\ohm} N-channel gate pull-down to GND                   & 0603   \\
\hline
\end{tabular}
\caption{High-side and low-side MOSFET switch components}
\label{tab:highside_parts}
\end{table}

\subsection{Logic Table}

\begin{table}[h]
\centering
\begin{tabular}{ccccl}
\hline
\textbf{GPIO A} & \textbf{NPN} & \textbf{AO3401A} & \textbf{GPIO B (AO3400A)} & \textbf{Load} \\
(high-side)     &              & (P-channel)      & (low-side)                &               \\
\hline
LOW  & OFF & OFF & LOW  & Off \\
LOW  & OFF & OFF & HIGH & Off \\
HIGH & ON  & ON  & LOW  & Off \\
HIGH & ON  & ON  & HIGH & \textbf{On} \\
\hline
\end{tabular}
\caption{High-side and low-side switch logic}
\label{tab:highside_logic}
\end{table}

\subsection{Example Code}

\begin{verbatim}
import RPi.GPIO as GPIO

HIGH_SIDE_PIN = 23   # controls AO3401A via MMBT3904
LOW_SIDE_PIN  = 24   # controls AO3400A directly

GPIO.setmode(GPIO.BCM)
GPIO.setup(HIGH_SIDE_PIN, GPIO.OUT, initial=GPIO.LOW)
GPIO.setup(LOW_SIDE_PIN,  GPIO.OUT, initial=GPIO.LOW)


def load_on() -> None:
    """Energise the load by enabling both high-side and low-side switches."""
    GPIO.output(HIGH_SIDE_PIN, GPIO.HIGH)
    GPIO.output(LOW_SIDE_PIN,  GPIO.HIGH)


def load_off() -> None:
    """De-energise the load. Low side is opened first."""
    GPIO.output(LOW_SIDE_PIN,  GPIO.LOW)
    GPIO.output(HIGH_SIDE_PIN, GPIO.LOW)
\end{verbatim}

\noindent The low-side switch is opened first in \texttt{load\_off()} to interrupt
current flow before the high side is released. This order avoids a brief condition
where the high side is off but the low side is still on, which would momentarily
short the load's negative terminal through the N-channel with no current path,
producing no harm but making the sequence logically cleaner.

	%	\section{Rotary Encoders}

\subsection{Overview}

Two rotary encoders are used in this design. One controls the alert buzzer volume and
the other controls the display brightness. Either encoder may be temporarily
reassigned by the application script to set other parameters as needed. The encoders
do not include push-button switches.

A rotary encoder is a digital position sensor. Unlike a potentiometer, it has no
end stops and produces no analog output. Instead, it produces a series of digital
pulses as the shaft is rotated, with the direction of rotation encoded in the phase
relationship between two output signals. Because the encoder has no absolute position,
the application maintains a software counter that is incremented or decremented on
each detent step.

\subsection{How a Rotary Encoder Works}

Inside the encoder are two switches, conventionally labelled A and B (sometimes CLK
and DT). The switches are arranged so that they open and close in a specific sequence
depending on the direction of rotation. As the shaft turns clockwise, A transitions
before B. As the shaft turns counter-clockwise, B transitions before A. The
relationship is a two-bit Gray code sequence:

\begin{table}[h]
\centering
\begin{tabular}{ccc}
\hline
\textbf{Step} & \textbf{A} & \textbf{B} \\
\hline
0 & 1 & 1 \\
1 & 0 & 1 \\
2 & 0 & 0 \\
3 & 1 & 0 \\
\hline
\end{tabular}
\caption{Gray code sequence for clockwise rotation. Counter-clockwise reverses the order.}
\label{tab:encoder_gray}
\end{table}

The simplest detection strategy is to trigger an interrupt on the falling edge of
signal A, then immediately read the state of signal B. If B is high at that moment,
rotation is clockwise; if B is low, rotation is counter-clockwise. This single-edge
approach is sufficient for a hand-turned control knob where speed is low.

\subsection{Contact Bounce and Debouncing}

The switches inside a mechanical rotary encoder bounce as they make and break contact,
producing multiple spurious transitions per detent step. Without debouncing, the
interrupt handler fires several times per click and the counter accumulates incorrect
values.

Two debouncing strategies are available:

\begin{itemize}
    \item \textbf{Hardware debouncing:} A small capacitor (typically \SI{10}{\nano\farad}
    to \SI{100}{\nano\farad}) from each encoder pin to ground, combined with a pull-up
    resistor (\SI{10}{\kilo\ohm}), forms an RC low-pass filter that smooths the bounce
    transients. This is the preferred approach for a production design as it requires no
    timing logic in software.

    \item \textbf{Software debouncing:} Record the timestamp of each interrupt and
    ignore subsequent triggers that arrive within a debounce window (typically
    \SI{5}{\milli\second} to \SI{50}{\milli\second}). The RPi.GPIO
    \texttt{add\_event\_detect()} call includes a built-in \texttt{bouncetime}
    parameter for exactly this purpose.
\end{itemize}

Both approaches may be combined. Hardware debouncing is recommended with a
\SI{100}{\nano\farad} ceramic capacitor on each encoder pin.

\subsection{Wiring}

Each encoder has two signal pins (A and B), a common pin, and optionally a supply
pin. The common is connected to ground. Each signal pin is pulled up to \SI{3.3}{\volt}
through a \SI{10}{\kilo\ohm} resistor (the Raspberry Pi internal pull-up may be used
instead, activated in software). A \SI{100}{\nano\farad} ceramic capacitor from each
signal pin to ground provides hardware debouncing.

\begin{verbatim}
  3.3V ---[ 10k ]---+--- Encoder A pin
                    |
                  100nF
                    |
                   GND
                   
  3.3V ---[ 10k ]---+--- Encoder B pin
                    |
                  100nF
                    |
                   GND

  Encoder Common --- GND
\end{verbatim}

\subsection{Pin Restrictions on the Raspberry Pi Zero}

All GPIO pins on the Raspberry Pi Zero support interrupts via the RPi.GPIO library.
There is no hardware restriction limiting interrupts to specific pins, unlike some
microcontrollers. The Linux kernel \texttt{gpio} subsystem services all GPIO interrupts
through a unified interface.

Practical limits to be aware of:

\begin{itemize}
    \item \textbf{Timing jitter:} Because the Raspberry Pi runs Linux, interrupt
    latency is not deterministic. Typical latency is \SI{1}{\milli\second} to
    \SI{5}{\milli\second} under light load, but may spike under heavy CPU or I/O
    activity. For a hand-turned knob this is entirely acceptable; for high-speed
    pulse counting it would not be.

    \item \textbf{Number of interrupts:} There is no fixed limit on the number of
    simultaneous GPIO interrupts. Each \texttt{add\_event\_detect()} call starts a
    kernel-level watch on that pin. In practice, a handful of interrupts is well
    within the capability of the system.

    \item \textbf{Power:} GPIO pins source or sink a maximum of \SI{16}{\milli\ampere}
    individually, with a total bank limit. Encoder signal pins draw negligible current
    when used as inputs with pull-ups; no power budget concern arises.

    \item \textbf{Voltage:} All GPIO inputs must remain within \SI{0}{\volt} to
    \SI{3.3}{\volt}. The encoders are passive switches to ground and are inherently
    safe provided the pull-ups are connected to \SI{3.3}{\volt} and not \SI{5}{\volt}.
\end{itemize}

\subsection{Setting Up Interrupts in Python}

The RPi.GPIO library provides \texttt{add\_event\_detect()} to register a callback
function that is invoked when a specified edge transition is detected on a GPIO pin.
The callback runs in a separate thread managed by the library, but in a simple
synchronous script the main loop can simply read shared variables that the callback
updates.

The general pattern for a single interrupt is:

\begin{verbatim}
import RPi.GPIO as GPIO

GPIO.setmode(GPIO.BCM)

PIN = 17
GPIO.setup(PIN, GPIO.IN, pull_up_down=GPIO.PUD_UP)

def on_edge(channel):
    # called automatically when the edge is detected
    pass

GPIO.add_event_detect(
    PIN,
    GPIO.FALLING,
    callback=on_edge,
    bouncetime=10        # ignore re-triggers within 10 ms
)
\end{verbatim}

\subsection{Two Encoders with Multiple Interrupts}

For two encoders, four GPIO pins are used (A and B for each encoder). An interrupt
is registered on the A pin of each encoder. When the interrupt fires, the handler
reads the B pin to determine direction and updates the appropriate software counter.

\begin{verbatim}
import RPi.GPIO as GPIO
import time

# --- Pin assignments ---
ENC1_A = 17    # Volume encoder
ENC1_B = 27
ENC2_A = 22    # Brightness encoder
ENC2_B = 23

# --- Shared state ---
volume     = 50    # 0-100
brightness = 50    # 0-100

GPIO.setmode(GPIO.BCM)

for pin in [ENC1_A, ENC1_B, ENC2_A, ENC2_B]:
    GPIO.setup(pin, GPIO.IN, pull_up_down=GPIO.PUD_UP)


def on_volume(channel):
    global volume
    direction = 1 if GPIO.input(ENC1_B) else -1
    volume = max(0, min(100, volume + direction))


def on_brightness(channel):
    global brightness
    direction = 1 if GPIO.input(ENC2_B) else -1
    brightness = max(0, min(100, brightness + direction))


GPIO.add_event_detect(ENC1_A, GPIO.FALLING,
                      callback=on_volume,     bouncetime=10)
GPIO.add_event_detect(ENC2_A, GPIO.FALLING,
                      callback=on_brightness, bouncetime=10)


def main():
    print("Turning encoders. Ctrl-C to exit.")
    try:
        while True:
            # Main loop runs normally; callbacks fire asynchronously
            print(f"Volume: {volume:3d}   Brightness: {brightness:3d}",
                  end="\r")
            time.sleep(0.1)
    except KeyboardInterrupt:
        pass
    finally:
        GPIO.cleanup()


if __name__ == "__main__":
    main()
\end{verbatim}

\noindent Key points in this implementation:

\begin{itemize}
    \item The B pin of each encoder is read inside the interrupt callback, not
    registered as a separate interrupt. Only the A pin triggers the interrupt.
    \item The \texttt{bouncetime=10} argument suppresses re-triggers within
    \SI{10}{\milli\second}, providing software debouncing in addition to any
    hardware RC filtering.
    \item The \texttt{global} keyword allows the callback (running in the RPi.GPIO
    interrupt thread) to modify the counter variables read by the main loop.
    \item The counters are clamped to their valid ranges inside each callback.
    \item \texttt{GPIO.cleanup()} is called on exit to release all pin reservations
    and remove event detection.
\end{itemize}

\subsection{Temporarily Reassigning an Encoder}

Because the encoder counters are simple integer variables, reassigning an encoder
to a different parameter requires only redirecting which variable the callback
modifies. The cleanest approach is to use a function that unregisters the existing
callback and registers a new one:

\begin{verbatim}
def reassign_encoder1(new_variable: list, label: str) -> None:
    """
    Temporarily redirect encoder 1 to control a different parameter.
    new_variable must be a mutable container, e.g. a list of one int.
    """
    GPIO.remove_event_detect(ENC1_A)

    def on_reassigned(channel):
        direction = 1 if GPIO.input(ENC1_B) else -1
        new_variable[0] = max(0, min(100, new_variable[0] + direction))

    GPIO.add_event_detect(ENC1_A, GPIO.FALLING,
                          callback=on_reassigned, bouncetime=10)
    print(f"Encoder 1 now controls: {label}")
\end{verbatim}

\noindent To restore the original volume function, call \texttt{reassign\_encoder1}
again with the \texttt{volume} container and label.

\subsection{Adafruit 1651 Touchscreen Interrupt (TSC2007)}

The Adafruit 1651 display uses the \textbf{TSC2007} resistive touch controller,
which communicates over I2C. The TSC2007 provides a \textbf{PENIRQ} pin (active
low) that is asserted whenever the screen is touched. Connecting this pin to a
GPIO and registering a falling-edge interrupt allows the application to respond to
touch events immediately, without polling the I2C bus in a continuous loop.

\begin{verbatim}
TOUCH_IRQ = 24    # PENIRQ from TSC2007

GPIO.setup(TOUCH_IRQ, GPIO.IN, pull_up_down=GPIO.PUD_UP)

def on_touch(channel):
    """Read touch coordinates from TSC2007 over I2C."""
    # Issue a read command to the TSC2007 over I2C here
    # e.g. using smbus2 or adafruit_tsc2007 library
    pass

GPIO.add_event_detect(TOUCH_IRQ, GPIO.FALLING,
                      callback=on_touch, bouncetime=50)
\end{verbatim}

\noindent The PENIRQ pin requires a pull-up to \SI{3.3}{\volt}; the internal
RPi pull-up activated by \texttt{GPIO.PUD\_UP} is sufficient. The PENIRQ line
is open-drain on the TSC2007 and safe to connect directly to the GPIO pin
provided the pull-up is to \SI{3.3}{\volt}.

\subsection{Parts List}

\begin{table}[h]
\centering
\begin{tabular}{llll}
\hline
\textbf{Reference} & \textbf{Description} & \textbf{Value / Part} & \textbf{Package} \\
\hline
ENC1, ENC2  & Rotary encoder, mechanical, no push button & e.g. Alps EC11  & Through-hole \\
R1--R4      & Signal pull-up resistors (A and B, each encoder) & \SI{10}{\kilo\ohm} & 0603 \\
C1--C4      & Hardware debounce capacitors (A and B, each encoder) & \SI{100}{\nano\farad} ceramic & 0603 \\
\hline
\end{tabular}
\caption{Rotary encoder components}
\label{tab:encoder_parts}
\end{table}

	
	
	\section{Magnetic Buzzer}
	
	\subsection{Overview}
	
	The alert sounder is a passive magnetic buzzer driven by a pulse-width modulated (PWM)
	signal from the Raspberry Pi Zero. A passive magnetic buzzer is preferred over an active
	buzzer or piezoelectric transducer because its tone is warm and rounded rather than harsh,
	and because both pitch and volume are controllable in software via the PWM frequency and
	duty cycle respectively.
	
	The buzzer operates from the \SI{5}{\volt} supply rail and draws approximately
	\SI{30}{\milli\ampere} to \SI{40}{\milli\ampere} at full drive. Since this exceeds the
	safe output current of a GPIO pin, and because the supply voltage exceeds the
	\SI{3.3}{\volt} GPIO logic level, a small-signal NPN bipolar junction transistor is
	used as a low-side switch between the buzzer and ground.
	
	\subsection{Circuit Description}
	
	The GPIO pin drives the base of the NPN transistor through a current-limiting resistor
	R1. When the GPIO is high, the transistor saturates and connects the buzzer's negative
	terminal to ground, completing the circuit through the \SI{5}{\volt} supply. When the
	GPIO is low, the transistor is cut off and the buzzer is de-energised.
	
	Because the buzzer is an inductive load, a flyback diode is placed across its terminals
	to suppress the reverse voltage spike that occurs when the transistor switches off.
	Without this diode the spike could damage the transistor.
	
	Volume is controlled by adjusting the PWM duty cycle. A duty cycle of approximately
	50\% produces full volume; reducing the duty cycle lowers the average power delivered
	to the buzzer and therefore reduces loudness. Pitch is controlled by adjusting the PWM
	frequency. The range of \SI{2000}{\hertz} to \SI{2500}{\hertz} produces a clear,
	pleasant tone that is audible in both laboratory and outdoor environments. The buzzer
	is rated at approximately \SI{85}{\decibel} to \SI{90}{\decibel} at \SI{10}{\centi\meter}.
	
	The complete circuit is shown below.
	
	\begin{verbatim}
		RPi GPIO (PWM) ---[ R1 1k ]---+--- MMBT3904 Base
		|
		MMBT3904 Emitter --- GND
		
		+5V ---[ Buzzer CEM-1203(42) ]---+--- MMBT3904 Collector
		|
		1N4148W (flyback)
		Anode  -> Collector side
		Cathode -> +5V side
	\end{verbatim}
	
	\subsection{Parts List}
	
	\begin{table}[h]
		\centering
		\begin{tabular}{llll}
			\hline
			\textbf{Reference} & \textbf{Part Number} & \textbf{Description} & \textbf{Package} \\
			\hline
			BZ1 & CEM-1203(42)  & Passive magnetic buzzer, \SI{5}{\volt}, \SI{2048}{\hertz} rated & 12\,mm through-hole \\
			D1  & 1N4148W       & Fast switching diode, flyback protection                        & SOD-123 \\
			Q1  & MMBT3904      & NPN small-signal BJT, \SI{40}{\volt}, \SI{200}{\milli\ampere}  & SOT-23  \\
			R1  & R1K-0805      & Resistor \SI{1}{\kilo\ohm}, base current limiter               & 0805    \\
			\hline
		\end{tabular}
		\caption{Magnetic buzzer circuit components}
		\label{tab:buzzer_parts}
	\end{table}
	
	\subsection{Programming Notes}
	
	The buzzer is controlled using the RPi.GPIO software PWM facility. Any GPIO pin may
	be used, but a hardware PWM pin (GPIO 18 or GPIO 12 on the Raspberry Pi Zero) is
	preferred because software PWM can exhibit jitter under system load, which may cause
	audible roughness in the tone.
	
	Key relationships are:
	
	\begin{itemize}
		\item \textbf{Frequency} sets pitch. The range \SI{1500}{\hertz} to
		\SI{3000}{\hertz} is recommended. Values below \SI{1000}{\hertz} may be too
		quiet for some magnetic buzzers, and values above \SI{4000}{\hertz} become
		increasingly harsh.
		\item \textbf{Duty cycle} sets volume. A value of 50 (50 percent) gives maximum loudness.
		Values below approximately 10 may be too quiet to be useful. Values above 50
		do not increase loudness and may introduce distortion.
		\item \textbf{Duration} is implemented as a blocking \texttt{time.sleep()} call.
		Interruption logic is not implemented in this version.
	\end{itemize}
	
	The GPIO pin must be initialised as an output before the PWM object is created. The
	PWM object should be started with a duty cycle of zero and frequency set before each
	call to avoid a brief click at startup.
	
	\subsection{Example Code}
	
	\begin{verbatim}
		import RPi.GPIO as GPIO
		import time
		
		BUZZER_PIN = 18  # Hardware PWM pin recommended
		
		GPIO.setmode(GPIO.BCM)
		GPIO.setup(BUZZER_PIN, GPIO.OUT)
		
		_pwm = GPIO.PWM(BUZZER_PIN, 2000)
		_pwm.start(0)
		
		
		def sound_alert(frequency: int = 2000,
		volume: int = 50,
		duration: float = 1.0) -> None:
		"""
		Sound the magnetic buzzer for a fixed duration.
		
		Parameters
		----------
		frequency : int
		Tone frequency in Hz. Recommended range 1500-3000.
		Default 2000 Hz.
		volume : int
		Duty cycle as a percentage, 0-50. Values above 50
		do not increase loudness. Default 50 (full volume).
		duration : float
		Duration in seconds to sound the alert. Blocking.
		Default 1.0 second.
		"""
		volume = max(0, min(50, volume))       # clamp to 0-50
		frequency = max(100, min(10000, frequency))  # clamp to safe range
		
		_pwm.ChangeFrequency(frequency)
		_pwm.ChangeDutyCycle(volume)
		time.sleep(duration)
		_pwm.ChangeDutyCycle(0)                # silence buzzer, leave PWM running
	\end{verbatim}
	
	\noindent To sound a two-second alert at \SI{2500}{\hertz} at three-quarter volume:
	
	\begin{verbatim}
		sound_alert(frequency=2500, volume=37, duration=2.0)
	\end{verbatim}
	
	\noindent To clean up GPIO resources on program exit:
	
	\begin{verbatim}
		_pwm.stop()
		GPIO.cleanup()
	\end{verbatim}
	
	
	
	\section{Extending MOSFETs with Shift Register}
	
	\subsection{Overview}
	
	The Raspberry Pi Zero provides a limited number of GPIO pins. Where multiple independent
	loads must be switched, a serial-to-parallel shift register allows eight MOSFET gates to
	be controlled from three GPIO pins. The device chosen is the \textbf{74HC595}, an
	8-bit serial-in, parallel-out shift register that operates from \SI{3.3}{\volt} and is
	widely available in an SMD package.
	
	Data is clocked into the shift register serially and latched to the eight parallel outputs
	simultaneously, ensuring that all outputs update at the same instant. Although the
	74HC595 can be driven by bit-banging three GPIO pins, the implementation here uses the
	hardware \textbf{SPI1} peripheral on the Raspberry Pi Zero for reliable timing and
	minimal CPU overhead. Only one MOSFET is energised at a time; the controlling byte
	therefore always has exactly one bit set.
	
	\subsection{Circuit Description}
	
	The Raspberry Pi Zero SPI1 peripheral provides three signals: MOSI (data), SCLK (clock),
	and CE (chip enable, used here as the latch signal). Each signal passes through a
	\SI{470}{\ohm} series resistor before reaching the 74HC595 to protect the GPIO pins
	from transient currents.
	
	The 74HC595 output enable pin ($\overline{\text{OE}}$) is tied permanently to ground so
	that the parallel outputs are always active. The master reset pin ($\overline{\text{MR}}$)
	is tied permanently to \SI{3.3}{\volt} to prevent accidental clearing of the register.
	A \SI{100}{\nano\farad} bypass capacitor is placed directly across the Vcc and GND
	pins of the IC.
	
	Each of the eight parallel outputs drives a MOSFET gate through a \SI{500}{\ohm}
	series resistor. A \SI{10}{\kilo\ohm} resistor from gate to ground ensures the gate
	is pulled firmly low when the corresponding output is de-asserted, preventing the
	MOSFET from floating into an undefined state.
	
	\subsection{Circuit Diagram}
	
	\begin{verbatim}
		Raspberry Pi Zero                  74HC595 (SOIC-16)
		(SPI1)                             +----------------+
		3.3V -----| Vcc(16) (8)GND |------ GND
		|                |
		GPIO 20 (MOSI) --[ 470R ]--------| DS   (pin 14)  |   [SER]
		GPIO 21 (SCLK) --[ 470R ]--------| SHCP (pin 11)  |   [SRCLK]
		GPIO 18 (CE0)  --[ 470R ]--------| STCP (pin 12)  |   [RCLK]
		|                |
		GND ----| /OE  (pin 13)  |  (always enabled)
		3.3V ----| /MR  (pin 10)  |  (master reset off)[SRCLR]
		|                |
		|   Q0  (pin 15) |--+
		|   Q1  (pin 1)  |--+  (Q1..Q7 identical)
		|   Q2  (pin 2)  |--+
		|   Q3  (pin 3)  |--+
		|   Q4  (pin 4)  |--+
		|   Q5  (pin 5)  |--+
		|   Q6  (pin 6)  |--+
		|   Q7  (pin 7)  |--+
		|       (pin 9)  |--x  (Q7S serial data out)
		100nF     |                |
		+----------------+
		--||-- Vcc to GND, close to IC
		
		One MOSFET channel (repeated x8):
		
		74HC595 Qn ---[ 500R ]---+--- MOSFET Gate (AO3400A or similar)
		|
		10k
		|
		GND                MOSFET Source --- GND
		MOSFET Drain  --- Load (-)
	\end{verbatim}
	
	\subsection{Parts List}
	
	\begin{table}[h]
		\centering
		\begin{tabular}{llll}
			\hline
			\textbf{Reference} & \textbf{Part Number} & \textbf{Description} & \textbf{Package} \\
			\hline
			U1      & 74HC595           & 8-bit shift register, serial-in parallel-out  & SOIC-16 \\
			Q1--Q8  & AO3400A           & N-channel MOSFET, \SI{30}{\volt}, \SI{5.7}{\ampere} & SOT-23 \\
			R1--R3  & R470-0603         & \SI{470}{\ohm} GPIO series resistors          & 0603 \\
			R4--R11 & R500-0603         & \SI{500}{\ohm} MOSFET gate resistors          & 0603 \\
			R12--R19 & R10K-0603        & \SI{10}{\kilo\ohm} gate pull-down resistors   & 0603 \\
			C1      & 100nF ceramic     & Bypass capacitor, Vcc to GND                  & 0603 \\
			\hline
		\end{tabular}
		\caption{Shift register MOSFET controller components}
		\label{tab:shiftreg_parts}
	\end{table}
	
	\subsection{Programming Notes}
	
	The SPI1 peripheral on the Raspberry Pi Zero is accessed via the \texttt{spidev} library.
	SPI mode 0 (CPOL=0, CPHA=0) is compatible with the 74HC595. A clock speed of
	\SI{500}{\kilo\hertz} is more than adequate; the 74HC595 is rated to several megahertz
	but there is no benefit to running faster for low-speed MOSFET switching.
	
	A single byte is transferred per operation. The byte has exactly one bit set to select
	one MOSFET channel, or is zero to de-energise all channels. The \texttt{xfer2()} call
	clocks out the data and automatically toggles the chip enable line, which is connected
	to the latch (STCP) pin, causing the outputs to update at the end of the transfer.
	
	Note that SPI1 on the Raspberry Pi Zero requires the overlay \texttt{spi1-1cs} to be
	enabled in \texttt{/boot/config.txt}:
	
	\begin{verbatim}
		dtoverlay=spi1-1cs
	\end{verbatim}
	
	\subsection{Example Code}
	
	\begin{verbatim}
		import spidev
		
		SPI_BUS = 1       # SPI1
		SPI_DEVICE = 0    # CE0
		
		spi = spidev.SpiDev()
		spi.open(SPI_BUS, SPI_DEVICE)
		spi.max_speed_hz = 500000
		spi.mode = 0
		
		
		def set_mosfet(channel: int) -> None:
		"""
		Energise a single MOSFET by channel number.
		
		Parameters
		----------
		channel : int
		Channel to activate, 0 through 7.
		Pass -1 or any out-of-range value to de-energise all channels.
		"""
		if 0 <= channel <= 7:
		byte = 1 << channel
		else:
		byte = 0
		spi.xfer2([byte])
		
		
		def all_off() -> None:
		"""De-energise all MOSFET channels."""
		spi.xfer2([0x00])
	\end{verbatim}
	
	\noindent To activate channel 4 then turn everything off:
	
	\begin{verbatim}
		set_mosfet(4)
		time.sleep(1.0)
		all_off()
	\end{verbatim}
	
	
	% ---------------------------------------------------------------------------
	\section{High-Side Control of 5\,V Load with MOSFET: Optional. Not used.}
	
	\subsection{Overview}
	
	In circuits requiring both high-side and low-side switching, the load is interrupted at
	both the positive supply terminal and the ground return. The low-side switch is an
	N-channel MOSFET (AO3400A), which is straightforward to drive from a \SI{3.3}{\volt}
	GPIO. The high-side switch requires a P-channel MOSFET because its source terminal
	connects to the positive supply rail.
	
	The P-channel MOSFET chosen is the \textbf{AO3401A}, the complementary P-channel
	counterpart to the AO3400A. Both devices share the SOT-23 package and are produced
	by the same manufacturer, making them a well-matched pair. The AO3401A is rated at
	\SI{-30}{\volt} Vds and \SI{-4}{\ampere} continuous drain current, well in excess of
	the anticipated load current.
	
	\subsection{The Level-Shifting Problem}
	
	A P-channel MOSFET turns on when its gate is pulled sufficiently below its source
	voltage. With the source at \SI{+5}{\volt}, the gate must be pulled close to ground
	to turn the device on, and must be pulled to \SI{+5}{\volt} to turn it off. A
	\SI{3.3}{\volt} GPIO cannot swing to \SI{+5}{\volt} and therefore cannot turn the
	device fully off. A small-signal NPN transistor is used to solve this: it pulls the
	gate to ground when the GPIO is high (turning the P-channel on), and allows the gate
	to be pulled to \SI{+5}{\volt} through a resistor when the GPIO is low (turning the
	P-channel off).
	
	\subsection{Circuit Description}
	
	The GPIO drives the base of an MMBT3904 NPN transistor through a \SI{10}{\kilo\ohm}
	base resistor. The collector of the NPN connects to the gate of the AO3401A. A
	\SI{10}{\kilo\ohm} resistor from the gate to \SI{+5}{\volt} pulls the gate high
	when the NPN is off, ensuring the P-channel is firmly off. When the NPN saturates,
	it pulls the gate to ground, turning the P-channel on.
	
	The logic is inverted by the NPN stage: GPIO high turns the high-side MOSFET on,
	and GPIO low turns it off. This matches the behaviour of the low-side N-channel
	switch, so both MOSFETs are enabled by a high logic level on their respective GPIO
	pins. The load is energised only when both GPIO pins are high simultaneously.
	
	When the supply voltage is \SI{5}{\volt} and the gate is pulled to ground by the NPN,
	Vgs of the AO3401A is \SI{0}{\volt} $-$ \SI{5}{\volt} $=$ \SI{-5}{\volt}. The
	AO3401A gate threshold is \SI{-1}{\volt}, so the device is driven well into
	conduction.
	
	\subsection{Circuit Diagram}
	
	\begin{verbatim}
		High-side control:
		
		RPi GPIO A --[ 10k ]--+-- MMBT3904 Base
		|
		MMBT3904 Emitter ----------- GND
		MMBT3904 Collector --+------ AO3401A Gate
		|
		10k
		|
		+5V
		
		+5V ----------- AO3401A Source
		AO3401A Drain ----+---- Load (+)
		|
		Load (-)
		|
		AO3400A Drain
		AO3400A Source --- GND
		AO3400A Gate --[ 500R ]--+-- RPi GPIO B
		|
		10k
		|
		GND
		
		Load energised when: GPIO A = HIGH and GPIO B = HIGH
	\end{verbatim}
	
	\subsection{Parts List}
	
	\begin{table}[h]
		\centering
		\begin{tabular}{llll}
			\hline
			\textbf{Reference} & \textbf{Part Number} & \textbf{Description} & \textbf{Package} \\
			\hline
			Q1 & AO3401A   & P-channel MOSFET, \SI{-30}{\volt}, \SI{-4}{\ampere}, high-side switch & SOT-23 \\
			Q2 & AO3400A   & N-channel MOSFET, \SI{30}{\volt}, \SI{5.7}{\ampere}, low-side switch  & SOT-23 \\
			Q3 & MMBT3904  & NPN BJT, \SI{40}{\volt}, \SI{200}{\milli\ampere}, gate driver         & SOT-23 \\
			R1 & R10K-0603 & \SI{10}{\kilo\ohm} NPN base resistor                                  & 0603   \\
			R2 & R10K-0603 & \SI{10}{\kilo\ohm} P-channel gate pull-up to \SI{5}{\volt}            & 0603   \\
			R3 & R500-0603 & \SI{500}{\ohm} N-channel gate series resistor                         & 0603   \\
			R4 & R10K-0603 & \SI{10}{\kilo\ohm} N-channel gate pull-down to GND                   & 0603   \\
			\hline
		\end{tabular}
		\caption{High-side and low-side MOSFET switch components}
		\label{tab:highside_parts}
	\end{table}
	
	\subsection{Logic Table}
	
	\begin{table}[h]
		\centering
		\begin{tabular}{ccccl}
			\hline
			\textbf{GPIO A} & \textbf{NPN} & \textbf{AO3401A} & \textbf{GPIO B (AO3400A)} & \textbf{Load} \\
			(high-side)     &              & (P-channel)      & (low-side)                &               \\
			\hline
			LOW  & OFF & OFF & LOW  & Off \\
			LOW  & OFF & OFF & HIGH & Off \\
			HIGH & ON  & ON  & LOW  & Off \\
			HIGH & ON  & ON  & HIGH & \textbf{On} \\
			\hline
		\end{tabular}
		\caption{High-side and low-side switch logic}
		\label{tab:highside_logic}
	\end{table}
	
	\subsection{Example Code}
	
	\begin{verbatim}
		import RPi.GPIO as GPIO
		
		HIGH_SIDE_PIN = 23   # controls AO3401A via MMBT3904
		LOW_SIDE_PIN  = 24   # controls AO3400A directly
		
		GPIO.setmode(GPIO.BCM)
		GPIO.setup(HIGH_SIDE_PIN, GPIO.OUT, initial=GPIO.LOW)
		GPIO.setup(LOW_SIDE_PIN,  GPIO.OUT, initial=GPIO.LOW)
		
		
		def load_on() -> None:
		"""Energise the load by enabling both high-side and low-side switches."""
		GPIO.output(HIGH_SIDE_PIN, GPIO.HIGH)
		GPIO.output(LOW_SIDE_PIN,  GPIO.HIGH)
		
		
		def load_off() -> None:
		"""De-energise the load. Low side is opened first."""
		GPIO.output(LOW_SIDE_PIN,  GPIO.LOW)
		GPIO.output(HIGH_SIDE_PIN, GPIO.LOW)
	\end{verbatim}
	
	\noindent The low-side switch is opened first in \texttt{load\_off()} to interrupt
	current flow before the high side is released. This order avoids a brief condition
	where the high side is off but the low side is still on, which would momentarily
	short the load's negative terminal through the N-channel with no current path,
	producing no harm but making the sequence logically cleaner.
	
	
	
	
	
\end{document}